\subsection{Inside the budget channel: resources against composition and needs}\label{inside-the-budget-channel-resources-against-composition-and-needs}

The resources-and-needs contribution is itself two things: the non-labour resources a household has, and the size and composition of the household those resources must cover. Separating them is not an arithmetic split of the joint cell. It requires two further counterfactual panels, each repriced through the tax-benefit system, and it is done here on the corrected partition of the budget fields into 58 resource fields, 46 composition and needs fields and 6 geographic fields, which travel with resources.

\Needspace*{12\baselineskip}\small\setlength{\tabcolsep}{3pt}\begin{longtable}[]{@{}
  >{\raggedright\arraybackslash}p{(\linewidth - 12\tabcolsep) * \real{0.1429}}
  >{\raggedright\arraybackslash}p{(\linewidth - 12\tabcolsep) * \real{0.1429}}
  >{\raggedright\arraybackslash}p{(\linewidth - 12\tabcolsep) * \real{0.1429}}
  >{\raggedright\arraybackslash}p{(\linewidth - 12\tabcolsep) * \real{0.1429}}
  >{\raggedright\arraybackslash}p{(\linewidth - 12\tabcolsep) * \real{0.1429}}
  >{\raggedright\arraybackslash}p{(\linewidth - 12\tabcolsep) * \real{0.1429}}
  >{\raggedright\arraybackslash}p{(\linewidth - 12\tabcolsep) * \real{0.1429}}@{}}
\caption{The subdivision of the resources-and-needs contribution into non-labour resources and household composition and needs, reported as ``resources / composition'' in every cell. Columns two and five are shares of that population's baseline inequality; the remaining columns are shares of the resources-and-needs channel itself, which the two cells divide. The couple cells come from two counterfactual panels repriced through the tax-benefit system on the corrected partition of 58 resource fields, 46 composition and needs fields and 6 geographic fields, and sum to the joint contribution to a residual of 1.4e-17 in index units. The single-adult cells come from panels priced on an earlier partition of the same fields, so they are reported for comparison and the two populations' channel shares are indicative rather than like for like.}\tabularnewline
\toprule\noalign{}
\begin{minipage}[b]{\linewidth}\raggedright
Index
\end{minipage} & \begin{minipage}[b]{\linewidth}\raggedright
Couples: share of baseline
\end{minipage} & \begin{minipage}[b]{\linewidth}\raggedright
Couples: share of channel
\end{minipage} & \begin{minipage}[b]{\linewidth}\raggedright
Couples: channel, equivalized
\end{minipage} & \begin{minipage}[b]{\linewidth}\raggedright
Singles: share of baseline
\end{minipage} & \begin{minipage}[b]{\linewidth}\raggedright
Singles: share of channel
\end{minipage} & \begin{minipage}[b]{\linewidth}\raggedright
Singles: channel, equivalized
\end{minipage} \\
\midrule\noalign{}
\endfirsthead
\toprule\noalign{}
\begin{minipage}[b]{\linewidth}\raggedright
Index
\end{minipage} & \begin{minipage}[b]{\linewidth}\raggedright
Couples: share of baseline
\end{minipage} & \begin{minipage}[b]{\linewidth}\raggedright
Couples: share of channel
\end{minipage} & \begin{minipage}[b]{\linewidth}\raggedright
Couples: channel, equivalized
\end{minipage} & \begin{minipage}[b]{\linewidth}\raggedright
Singles: share of baseline
\end{minipage} & \begin{minipage}[b]{\linewidth}\raggedright
Singles: share of channel
\end{minipage} & \begin{minipage}[b]{\linewidth}\raggedright
Singles: channel, equivalized
\end{minipage} \\
\midrule\noalign{}
\endhead
\bottomrule\noalign{}
\endlastfoot
Gini & 32.88 / 18.33 & 64.2 / 35.8 & 57.1 / 42.9 & 24.38 / 9.82 & 71.3 / 28.7 & 46.1 / 53.9 \\
Atkinson(1) & 49.14 / 11.46 & 81.1 / 18.9 & 73.7 / 26.3 & 34.38 / 12.52 & 73.3 / 26.7 & 48.3 / 51.7 \\
Atkinson(2) & 44.42 / 12.79 & 77.6 / 22.4 & 68.3 / 31.7 & 31.33 / 12.24 & 71.9 / 28.1 & 44.2 / 55.8 \\
GE(0) & 49.15 / 11.42 & 81.2 / 18.8 & 73.8 / 26.2 & 34.62 / 12.57 & 73.4 / 26.6 & 48.4 / 51.6 \\
GE(1) & 55.49 / 9.44 & 85.5 / 14.5 & 80.4 / 19.6 & 39.26 / 13.47 & 74.4 / 25.6 & 53.4 / 46.6 \\
GE(2) \(=CV^2/2\) & 63.15 / 7.01 & 90.0 / 10.0 & 87.7 / 12.3 & 47.04 / 15.81 & 74.8 / 25.2 & 58.7 / 41.3 \\
\end{longtable}

Three statements are supported, and a fourth that suggests itself is not.

First, \textbf{non-labour resources lead the channel in both populations on the household basis, under all six indices} --- 6 of six for couples and 6 of six for single adults. On the Gini the couples channel divides 64.2 to 35.8, which is 32.88 and 18.33 per cent of total inequality; the single-adult channel divides 71.3 to 28.7, or 24.38 and 9.82 per cent of total.

Second, \textbf{the resources share of the channel is widest under GE(2) and narrowest under the Gini in both populations.} The Gini is the index least sensitive to the tails, and it is where composition matters most; the squared coefficient of variation is the most tail-sensitive, and it is where resources dominate. The spread is large --- for couples, from 64.2 per cent of the channel under the Gini to 90.0 under GE(2) --- so the division of this channel is considerably more index-sensitive than the four-way decomposition above it.

Third, \textbf{equivalizing raises the composition share in both populations, under all 6 of the six indices.} This is not an artefact and it is not a surprise: the equivalence scale is owned by the composition operator, so a state in which composition is equalized is also a state in which every household is put on a common scale. Equalizing composition therefore removes both the direct effect of household size on the budget and the effect of size on the scale by which the resulting level is deflated. For single adults the effect is strong enough to reverse the ordering: on the equivalized basis resources lead the channel in only 2 of the six indices, and composition leads under the Gini, both Atkinson indices and GE(0). For couples the ordering survives equivalization, under all 6 of the six.

The statement that does \emph{not} survive is a comparison between the two populations. On the Gini, composition takes a larger share of the couples channel than of the single-adult channel, 35.8 against 28.7 per cent, which invites the reading that household composition matters more where there is a household to compose. That ordering holds under 1 of the six indices, the Gini, and reverses under the other five; on the equivalized basis it reverses under all six. We therefore do not report it as a finding. The two populations' panels also rest on different partitions of the budget fields, which is a second reason to treat the cross-population comparison of channel shares as indicative.

